\documentclass[12pt]{article}
\usepackage[utf8]{inputenc}
\usepackage[T1]{fontenc}
\usepackage{hyperref}
\title{Sango Language Resources for K-12 Curriculum Development}
\author{Compiled by Claude Sonnet 5.5 (claude-sonnet-5-5) for Andr\'as Kornai}
\date{}
\begin{document}
\maketitle

\section{Introduction}
Sango (ISO 639-3 sag; Glottocode sang1328) is a Ngbandi-based creole of the Central African Republic, widely used as a lingua franca. French is still dominant in official communication and schooling, while Sango is the main language of everyday speech. The searches found dictionary and orthography resources but no school textbooks.

\section{Primary Educational Resources}
No Sango school textbooks or national curriculum documents were found in the search results.
\begin{itemize}
\item Wikibooks open Sango course (a learner course, not a school curriculum; the page notes it needs more contributors): \url{https://en.wikibooks.org/wiki/Sango}
\end{itemize}

\section{Linguistic and Reference Materials}
\begin{itemize}
\item Online Sango dictionary on Webonary, with entries in the official orthography and French, English and German equivalents: \url{https://www.webonary.work/sango/overview/introduction/}
\item The Sango orthography (Webonary): \url{https://www.webonary.work/sango/?p=146823}
\item Sango alphabet page (22 letters, Webonary): \url{https://www.webonary.work/sango/overview/alphabet/?lang=en}
\item Sango-English / English-Sango Dictionary and Phrasebook by Marcel Diki-Kidiri and John Lechner, Hippocrene Books (about 4,000 entries; listings disagree on the publication date): \url{https://www.hippocrenebooks.com/store/p512/SangoDictionaryandPhrasebook.html}
\item Omniglot page on Sango and its writing: \url{https://www.omniglot.com/writing/sango.php}
\item Atlas of Pidgin and Creole Language Structures survey for Sango: \url{https://apics-online.info/surveys/59}
\end{itemize}

\section{Religious and Community Resources}
The searches indicate that Bible translations and a hymnal were published in the 1960s (a 1966 Bible and a 1968 hymnal are mentioned), but no URL for a scripture text was seen, so none is cited. The searches surfaced no community organisation pages.

\section{Digital and Audio Resources}
Nothing audio-specific was confirmed. A Radio Ndeke Luka listing was mentioned only as a link on a resources page, and no station URL was seen. The page below is an aggregator of unverified quality:
\begin{itemize}
\item Sango language resources list: \url{https://universeofmemory.com/sango-language-resources/}
\end{itemize}

\section{Curriculum Development Context}
Sango has an official orthography (Decree 84.025 of 1984) and a usable online dictionary, which gives a base for materials. School instruction, however, remains mainly in French, and no Sango-medium curriculum was found. Because Sango is mostly a second language, materials would need to serve both first- and second-language learners.

\end{document}